\scanpage{191}
% l. 01
( this \unsure{works} \unsure{without} a finiteness hypothesis on
% l. 02
$\Gamma$, the universal case being the one where $\Gamma$
% l. 03
is the Teichmüller \unsure{group} of $\pi'$ \unsure{fixing} $i$
% l. 04
i.e. $\operatorname{Autext}_{\mathrm{lac}}(\pi';i)\simeq\operatorname{Aut}_{\mathrm{lac}}(\pi)$ ). I
% l. 05
say that if $\Gamma$ operates \uline{faithfully} \uline{on $\pi$}
% l. 06
i.e. if its outer operation \uline{on}
% l. 07
\uline{$\pi'$} is faithful, and if $\Gamma$ is \uline{finite},
% l. 08
Then $\Gamma$ is necessarily cyclic \inserted{(case $\Gamma=\Gamma^0$)}
% l. 09
or dihedral ($\Gamma\ne\Gamma^0$), and that if in
% l. 10
addition $\Gamma=\Gamma^0$, then \deleted{\unsure{$\mathcal E_i'$ is $\simeq$}} \inserted{\unsure{$\mathcal E_i\simeq\mathbb Z$|$\mathcal E_i'\simeq\mathbb Z$} i.e.} $\exists!$ isomorphism
% l. 11
\unsure{$T\xrightarrow{\sim}\mathcal E_i$|$T\xrightarrow{\sim}\mathcal E_i'$} such that \unsure{$T\xrightarrow{k_i'}\mathcal E_i'$} identifies with
% l. 12
$n\operatorname{id}_T$ (so that we find \unsure{$\Gamma\simeq\mathcal E_i'/n\mathcal E_i'$}
% l. 13
$\simeq T\otimes_{\mathbb Z}\mathbb Z/n\mathbb Z$, if $n=\operatorname{card}\Gamma$ ...

% l. 14
Changing notation, this amounts to the
% l. 15
\uline{Theore}\textsuperscript{\unsure{m}} Let $\Gamma$ be a finite group operating
% l. 16
faithfully on an \uline{\unsure{outer|outer\textsuperscript{ly}}}
% l. 17
group with loops $\pi$ \inserted{of type \deleted{\illegible{1}} $(g,\nu+1)$ ($\nu\geq0$)}, leaving an $i\in I(\pi)$ fixed,
% l. 18
Then $\Gamma$ is cyclic (if $\Gamma=\Gamma^0$) or
% l. 19
dihedral (if $\Gamma\ne\Gamma^0$), and in the extension
% marge, en bas à gauche
\margin{\unsure{to the} same there is \uline{equivalence}, if in the theorem, we assume the type $(g,\nu)$ anabelian --- otherwise the theorem is a little more general}
\mspage{190}


\scanpage{192}
\mspage{96}
% l. 01
corresponding \deleted{\illegible{2}} \unsure{$\mathcal E$} of $\Gamma$ by $\pi$,
% l. 02
If \unsure{$\mathcal E_i$} is the normalizer of an $L_i$ in
% l. 03
\unsure{$\mathcal E$} [ \inserted{so that we have} \deleted{an} \unsure{$\mathcal E_i\to\Gamma$} epimorphic i.e. an exact sequence
% l. 04
\unsure{$1\to L_i\to\mathcal E_i\to\Gamma\to1$} ]. the inverse image \unsure{$\mathcal E_i^0$} of
% l. 05
$\Gamma^0$ is $\simeq\mathbb Z$ so (if $n=\operatorname{card}\Gamma^0$)\inserted{\unsure{$=[\mathcal E_i^0:L_i]$}}
% l. 06
\deleted{\unsure{$n\operatorname{id}$}} $x\to x^n$ is an isomorphism \deleted{and} \unsure{$\mathcal E_i^0\simeq L_i$}, which
% l. 07
taking $k_i$ into account gives an
% l. 08
isomorphism \unsure{$\mathcal E_i^0\simeq T$}, whose composite with
% l. 09
\deleted{\unsure{$\mathcal E_i$}} \unsure{$T\xrightarrow[\sim]{k_i}L_i\subset\mathcal E_i^0$} is $n\operatorname{id}_T$, so
% l. 10
that we have a canonical isomorphism
% l. 11
\[\unsure{\Gamma^0\simeq\mathcal E_i^0/L_i\simeq T/nT}\]
% l. 12
(obviously independent of the choice of $L_i$ ...)

% l. 13
\uline{Proof}. Let us consider the extension
% l. 14
\unsure{$\mathcal E_i^0$} of $\Gamma^0$ by $L_i\simeq T$ ($\overset{\text{\unsure{noncanonical}}}{\simeq}\mathbb Z$),
% l. 15
since \unsure{$\mathcal E_i^0$|$\Gamma^0$} operates trivial\textsuperscript{ly} on $T$ \inserted{torsion-free}, this
% l. 16
extension (by the cohomology exact sequence
% l. 17
associated with $0\to T\to T\otimes_{\mathbb Z}\mathbb Q\to T\otimes_{\mathbb Z}\mathbb Q/\mathbb Z\to0$)
% l. 18
is canonical\textsuperscript{ly} isomorphic to \deleted{\unsure{an}} the extension
% l. 19
defined by a homomorphism $\Gamma^0\to T_n=T/nT$,
% l. 20
as inverse image of the extension
% l. 21
$0\to T\xrightarrow{n}T\to T_n\to0$ of $T_n$ by $T$. I say that
\mspage{191}


\scanpage{193}
% l. 01
this homomorphism is an isomorphism (from which \unsure{would follow|will follow}
% l. 02
the other assertions), or what amounts
% l. 03
to the same since $\Gamma^0$ and $T_n$ are both
% l. 04
of order $n$, that it is a monomorphism.
% l. 05
Replacing $\Gamma^0$ by the kernel of $\Gamma^0\to\mathbb Z/n\mathbb Z$,
% l. 06
we may assume that the homomorphism \inserted{in question} is zero, and we must
% l. 07
prove that this implies that the action
% l. 08
of $\Gamma^0$ is trivial (which, since by hypothesis
% l. 09
the action of $\Gamma$ is faithful, implies $\Gamma^0=\{1\}$,
% l. 10
OK.) So \unsure{theorem} is reduced to the

% l. 11
\uline{fundamental lemma}. Every \inserted{direct} outer automorphism
% l. 12
$u$ of finite order \inserted{\unsure{$n$}} of a group with
% l. 13
loops $\pi$, which fixes a loop class
% l. 14
$i$ and is such that the extension of $\mathbb Z/n\mathbb Z$
% l. 15
by $L_i\in i$ defined by $u$ is trivial,
% l. 16
is trivial.

% l. 17
The hypothesis means that $u$ lifts
% l. 18
to an automorphism $u_0$ of $\pi$, which
% l. 19
normalizes $L_i$, and which is also of order \unsure{$n$},
% l. 20
(or of finite order, this amounts to the \unsure{same|$\Rightarrow$} taking into account
% l. 21
that \deleted{\unsure{the elements of}} $T$ has no elements of finite order)
\mspage{192}


\scanpage{194}
\mspage{97}
% l. 01
then $u_0$ is trivial, i.e. this is equivalent to the
% l. 02
\uline{Corollary}. Every \inserted{direct} automorphism of finite order \inserted{$u_0$} of
% l. 03
a group with loops $\pi$, \inserted{which is direct and} which normalizes
% l. 04
a loop subgroup $L_i$ (i.e. which
% l. 05
centralizes $L_i$), is reduced to the identity.
% marge supérieure gauche
\margin{In the lemma or its \unsure{corollary}}
\margin{\uline{N.B.} (the case $g=0$, $\nu+1=1$ or $\nu+1=2$ is trivial, the case $\nu+1=3$ ($(g,\nu)=(0,2)$ abelian!) is not trivial \unsure{on the other hand}, nor the case $g=1$, $\nu+1=1$ i.e. $(g,\nu)=(1,0)$ abelian. Yet the result should be valid also in \unsure{this|these} cases.)}

% l. 06
To prove it, we are reduced \inserted{at once} to the
% l. 07
case where $u_0$ is \deleted{of prime order $p$} \inserted{such that $u_0^p=\operatorname{id}$, where}
% l. 08
$p$ is a prime \unsure{number} --- i.e. $u_0$ corres-
% l. 09
ponds to an operation in the strict sense
% l. 10
(not on\textsuperscript{ly} outer) of $\mathbb Z/p\mathbb Z$ \deleted{\unsure{by}} \unsure{by automorphisms}
% l. 11
of $\pi$. But (whether or not we can
% l. 12
find such a $p$) let us consider the case where
% l. 13
we know that the \uline{outer} operation of
% l. 14
$\Gamma=\mathbb Z/n\mathbb Z$ [$n=$ order of $u$] on $\pi$ is realized
% l. 15
geometrical\textsuperscript{ly}, by a (faithful) operation
% l. 16
of $\Gamma=\mathbb Z/n\mathbb Z$ on a $U$ of type $(g,\nu+1)$.
% l. 17
\deleted{\unsure{$\pi$-conjugacy classes}, \illegible{1}}
% marge inférieure gauche
\margin{$U=X\setminus S$, $S\simeq I$, $X$ compact of genus $g$ \unsure{the puncture|the part} $s_i$ of $S$ corresponding to a fixed point of $\Gamma$ operating on $U$.}

% l. 18
\deleted{We know that the (\unsure{liftings}) of the outer automorphisms
% l. 19
to \deleted{\unsure{automorphisms}} \inserted{a genuine operation},
% l. 20
correspond to the fixed points \illegible{1} of $\Gamma$
% l. 21
in $U$; \deleted{and the extension of} \inserted{let $a$ be the fixed point of $\Gamma$}
% l. 22
corresponding to the given operation, \unsure{so}
% l. 23
that we have \deleted{\unsure{a}} canonical isomorphism \unsure{compatible with $\Gamma$}
% l. 24
$\pi=\pi_1(U,a)$.}
\mspage{193}


\scanpage{195}
% l. 01
We express the hypothesis of the \inserted{\unsure{``genuine''}} operation of $\Gamma$ on
% l. 02
$\pi$, \deleted{\illegible{1}} \unsure{centralizing} an $L_i$, by saying that
% l. 03
the outer operation gives a trivial extension
% l. 04
of $\Gamma$ by $L_i\simeq T$. But we know
% l. 05
moreover in the geometric case \inserted{and faithful operation} that it
% l. 06
is \uline{not} trivial !

% l. 07
It would therefore suffice, \deleted{\illegible{1}} to prove the
% l. 08
fundamental lemma, to know that every
% l. 09
outer action \inserted{[faithful]} of a cyclic group
% l. 10
on a group with loops of type \unsure{$(g,\nu+1)$}
% l. 11
is realizable, and it even suffices to know this
% l. 12
for a cyclic group of
% l. 13
prime order. \inserted{\unsure{is}} \unsure{Now} unless I am mistaken this result
% l. 14
is \unsure{well} known (even for \unsure{solv-}
% l. 15
able groups ...). --- so that the funda-
% l. 16
mental lemma seems proved. I have only
% l. 17
a doubt whether it is proved in the
% l. 18
general case of a $(g,\nu)$, \inserted{\deleted{\illegible{1} $g\geq2$}} or only
% l. 19
for \inserted{$g\geq2$,} $\nu=0$. But \deleted{in truth, \illegible{2}, \inserted{this suffices}}
% l. 20
\deleted{\illegible{2} \inserted{for $g\geq2$}}
% marge inférieure gauche
\margin{as a theorem \unsure{of existence} \unsure{of points} fixed by a group \illegible{1} \unsure{operating} \unsure{of} Teichmüller}

% l. 21
\deleted{\uline{Lemma} Let \illegible{1} be an outer operation}
\mspage{194}


\scanpage{196}
\mspage{98}
% l. 01
\deleted{of a finite group $\Gamma$ on a \illegible{1}} if this
% l. 02
is indeed so, I think that (for $g\geq2$ at the very
% l. 03
least) we get by with dévissage
% l. 04
for the case of arbitrary $\nu$ --- and the case $g=0,1$
% l. 05
would require \inserted{\unsure{also}} a separate treatment.
% l. 06
I will come back to this later, and
% l. 07
prefer for the moment to admit the ``fundamental
% l. 08
lemma'', and examine consequences
% l. 09
and corollaries of it.

% l. 10
For an outer \inserted{faithful} action of a finite $\Gamma$
% l. 11
on an \uline{anabelian} group with loops $\pi$, \unsure{we}
% l. 12
\deleted{\unsure{have established}} corresponding to an extension \unsure{$\mathcal E$} of $\Gamma$ by $\pi$,
% l. 13
we have therefore established a) that the partial
% l. 14
splittings of it can only be made
% l. 15
over \deleted{\unsure{cyclic}} subgroups $\Gamma'$ of $\Gamma$
% l. 16
such that $\Gamma'{}^0$ \unsure{is} cyclic, and $\Gamma'$ dihedral
% l. 17
if $\Gamma'\ne\Gamma'{}^0$, and that b) for every class
% l. 18
of \illegible{2} of such partial splittings
% l. 19
\unsure{over} $\Gamma'\subset\Gamma^0$, we have a corresponding canonical
% l. 20
isomorphism $\Gamma'\simeq T_n$ ($=T/nT$), where $n=\operatorname{card}\Gamma'$.
% marge inférieure gauche
\margin{i.e. we have established the existence of a canonical \unsure{map} \illegible{1} from the \unsure{elements} of finite order of $\operatorname{Aut}_{\mathrm{lac}}(\pi)$ into \unsure{$T\otimes\mathbb Q/T$}, satisfying \unsure{the} conditions \unsure{examined} previously}

% l. 21
These are results which we had
\mspage{195}


\scanpage{197}
% l. 01
previously obtained for the case of a
% l. 02
\uline{realizable} operation. [There was moreover no reason
% l. 03
to restrict ourselves to the anabelian case,
% l. 04
as long as we \unsure{assume} $\pi\ne0$ \inserted{(essentially empty case!)},
% l. 05
which includes the \inserted{abelian} cases $g=0$, $\nu=2$ and
% l. 06
$g=1$, $\nu=0$, for which a direct treatment
% l. 07
is possible, and has already been given
% l. 08
essentially, those cases \deleted{\unsure{being} realizable,} \inserted{\unsure{where|if one}}
% l. 09
starts with an \uline{extension} (not with an
% l. 10
``outer'' action i.e. here an action plain and simple
% l. 11
of $\Gamma$ on $\pi\simeq\mathbb Z$ or $\mathbb Z^2$), \unsure{being} always \unsure{realizable}.
% l. 12
(\uline{N.B.} \deleted{\unsure{the hypothesis}} \inserted{In this} case, the hypothesis of a
% l. 13
faithful action is replaced by the one that
% l. 14
the extension is \deleted{\unsure{not splittable}} \inserted{a \uline{product} extension} \unsure{over|under} no
% l. 15
subgroup $\Gamma'\subset\Gamma^0$ with $\Gamma'\ne1$ ....).

% l. 16
There remain however other results of the
% l. 17
realizable case, which would have to be examined
% l. 18
in the general case :

% l. 19
c) If $\Gamma$ is a \deleted{\unsure{discrete}} finite subgroup
% l. 20
of $\operatorname{Aut}_{\mathrm{lac}}(\pi)$ i.e. a finite group operating \uline{\unsure{faithfully}}
% l. 21
on $\pi$, do we have $\pi^\Gamma=\{1\}$ ?
% marge inférieure gauche
\margin{We may assume $\Gamma=\Gamma^0$ hence $\Gamma$ generated by a direct automorphism \unsure{$u$} or $\Gamma$ generated by an \unsure{anti-automorphism} of order 2}
\mspage{196}


\scanpage{198}
\mspage{99}
% l. 01
d) \unsure{If} $\Gamma',\Gamma''$ are finite subgroups of \unsure{$\mathcal E$}
% l. 02
\unsure{(extension} of the finite group $\Gamma$ by $\pi$, corresponding
% l. 03
to an \deleted{\unsure{respect}} outer \inserted{(preserving the orientation)} faithful operation), \deleted{\illegible{1}}
% l. 04
such that $\Gamma'\cap\Gamma''\ne(1)$, then \deleted{does one of
% l. 05
the two contain the other} ?
\inserted{are $\Gamma'$ and $\Gamma''$ contained in the $\pi$-conjugate of a finite subgroup $\Gamma'''$ of \unsure{$\mathcal E$} [i.e. every section sub\unsure{group} $\ne1$ is contained in a unique \unsure{maximal} section subgroup, modulo $\pi$-conjugation}
% marge supérieure gauche
\margin{\uuline{N.B.} Would become false by \deleted{\illegible{1}} \unsure{placing} in the group $\operatorname{Aut}_{\mathrm{lac}}(\pi)$ \illegible{2} \unsure{one as a subgroup} \illegible{2} \unsure{one of the two subgroups} \unsure{contained in} \illegible{4}}

% l. 06
e) For every \unsure{subgroup}, \illegible{2} \deleted{\illegible{1}} $\Gamma'$ of $\Gamma$,
% l. 07
is \unsure{the set} of classes \deleted{\unsure{of extensions}} of $\pi$-conjugation
% l. 08
of liftings of $\Gamma'$ into \unsure{$\mathcal E$} finite ?
% l. 09
[We may assume $\Gamma'=\Gamma$, and $\Gamma$ cyclic (if $\Gamma=\Gamma^0$) or dihedral
% l. 10
(otherwise)]
% l. 11
\inserted{In case \unsure{c)|e)}, we may assume $\Gamma$ cyclic of prime order, and
% l. 12
it is OK if it is established that an \unsure{outer operation} \illegible{3}
% l. 13
\unsure{that} \illegible{4} \unsure{$\pi_{g,\nu}$} \deleted{\illegible{3}} \unsure{d).}} \deleted{By applying}
% marge gauche, au milieu
\margin{\unsure{In} the same way e) is established if \unsure{we} admit that the outer operations of \unsure{cyclic groups} \unsure{on a} $\pi\simeq\pi_{g,\nu}$ are realizable}

% l. 14
\deleted{Example. let us examine \illegible{2}
% l. 15
the \unsure{juncture|junction} ``drilling holes'', $\Gamma'$ and $\Gamma''$
% l. 16
operating by outer automorphisms on a $(\pi';i)$
% l. 17
$i\in I(\pi')$ --- their action on $\pi$ is
% l. 18
deduced from that on $(\pi',i)$ by ``plugging
% l. 19
the hole $i$''. That said, $\Gamma'$ and $\Gamma''$ are identified
% l. 20
with the finite subgroups of $\operatorname{Autext}_{\mathrm{lac}}(\pi';i)$
% l. 21
We have better :
% l. 22
\unsure{$\Gamma=\Gamma'\cap\Gamma''\ne(1)$}}

% Un trait oblique part du coin inférieur gauche du cadre et traverse « Soit » et « et considérons » ; sa portée comme biffure est incertaine.
\unsure{\deleted{
% l. 23
\uline{Theorem}. Let $\Gamma_0$ \deleted{\illegible{1}} be a finite subgroup of
% l. 24
$A=\operatorname{Aut}_{\mathrm{lac}}(\pi)^0$, \inserted{$\Gamma_0\ne1$}, \deleted{\illegible{1} and let us consider} Then the normalizer $N$
% l. 25
of \unsure{$\Gamma$|$\Gamma_0$} in $A$ is equal to its centralizer \unsure{$Z$|$\mathcal Z$}
}}
\margin{\uuline{Ultra strong!} \unsure{If $g=0$}}
\mspage{197}


\scanpage{199}
% Un trait oblique traverse le corps de cette page ; sa portée comme biffure est incertaine.
\unsure{\deleted{
% l. 01
and the set of elements of finite order
% l. 02
of \unsure{$Z$|$\mathcal Z$} is a cyclic group \inserted{\unsure{$F$|$\Gamma$}}, \deleted{contained
% l. 03
in} which is also the largest finite
% l. 04
subgroup of $A$ containing $\Gamma_0$, and is
% l. 05
contained in the center of \unsure{$Z$|$\mathcal Z$}.

% l. 06
The fact that the normalizer $N$
% l. 07
centralizes $\Gamma_0$ comes from the fact that we
% l. 08
have a canonical isomorphism (compatible with the \unsure{transports}
% l. 09
of \unsure{structures})
% l. 10
\[\Gamma_0\xrightarrow{\sim}T_{n_0}\qquad(\text{if }n_0=\operatorname{card}\Gamma_0)\]
% l. 11
(N.B. if $g\in\operatorname{Aut}_{\mathrm{lac}}(\pi)$ reversing the orientation,
% l. 12
and \unsure{normalizes|centralizes} $\Gamma_0$, then \deleted{\unsure{$g$}} $(\operatorname{int}g)$ induces
% l. 13
the symmetry $\sigma\to\sigma^{-1}$ on $\Gamma_0$, by the same
% l. 14
argument) \deleted{\illegible{1}} Let us note that if
% l. 15
$u$ is an element of finite order of \unsure{$Z$|$\mathcal Z$}, so
% l. 16
that $\Gamma_0,u$ generate a
% l. 17
finite commutative subgroup of \unsure{$Z$|$\mathcal Z$}, \deleted{we see that} \inserted{as we know}
% l. 18
that a finite subgroup of $A$ is cyclic,
% l. 19
it follows that $u$ is contained in
% l. 20
a cyclic subgroup of \unsure{$Z$|$\mathcal Z$} \unsure{contain\textsuperscript{ing}} $\Gamma_0$,
% l. 21
\deleted{\unsure{therefore}} say $\Gamma'$. let us now choose
% l. 22
a geometric realization of the \inserted{outer} action of
% l. 23
$\Gamma_0$ on \deleted{\illegible{1}} $\pi$, say $(X,S)$ with $U=X-S$ ;
}}
\mspage{198}


\scanpage{200}
\mspage{100}
% Un long trait oblique traverse le texte ; sa portée comme biffure est incertaine.
\unsure{\deleted{
% l. 01
we know that the lifting of the outer
% l. 02
action of $\Gamma_0$ to an action is
% l. 03
described by a unique point $a\in U$.
% l. 04
\deleted{Let} I say that $\Gamma'$ \deleted{is realizable} (or its
% l. 05
\deleted{outer} action on $\pi\simeq\pi_1(U,a)$) is realizable
% l. 06
in a \uuline{unique way} by an action of $\Gamma'$
\margin{\uuline{not true!}}
\inserted{and inducing the already given action of \unsure{$\Gamma$|$\Gamma_0$}}
% l. 07
on $(X,S)$), fixing $a$, the fact that it
% l. 08
is \inserted{indeed} realizable in this way comes from the fact
% l. 09
(which we admit) that every outer
% l. 10
action of a cyclic group \inserted{(here $\Gamma'$)} on a
% l. 11
group with loops \inserted{$\pi$} of type $(g,\nu)$ is realizable,
% l. 12
and from the uniqueness of such a realization
% l. 13
(in the case of $\Gamma_0$). For \uline{uniqueness},
% l. 14
this \deleted{\illegible{1}} comes down to the

% l. 15
\uline{Lemma}. Let $u_0$ be an automorphism of finite order $n_0\ne1$
% l. 16
of an \inserted{orientable} surface \deleted{\unsure{assumed} $X^{u_0}\ne\emptyset$.} \inserted{$U=X-S$, $X$ (compact, \deleted{\illegible{1}}, and} leaving fixed \inserted{a point $a\in U$} \illegible{1}
% l. 17
let $d\in\mathbb N^*$, (If there exists an automorphism \unsure{$u$} of
% l. 18
\deleted{\illegible{1}} \inserted{$(X,S)$} such that $u^d=u_0$, \inserted{$u(a)=a$}, then $u$ is uniquely
\margin{\uline{false!}}
% l. 19
determined by its \deleted{\unsure{class of}} action on $\pi_1(U,a)$;
% l. 20
and in any event the cyclic group
% l. 21
of automorphisms of $(X,S)$ generated by $u$ is \unsure{unique\textsuperscript{ly}}
% l. 22
\inserted{determined (it is of order $nd$).}
}}
\mspage{199}
